\documentclass[12pt]{article}
\usepackage[T1]{fontenc}
\usepackage[utf8]{inputenc}
\usepackage{lmodern}
\usepackage{textcomp}
\usepackage{enumitem}
\usepackage{geometry}
\geometry{margin=1in}
\begin{document}
\begin{center}
Answer Key - Mathematics - Graduate Level Assessment 23 \
\small (Answers are ordered exactly as in the question paper.)
\end{center}

\section*{Multiple Choice}
\begin{enumerate}[label=\arabic*.]
\item \textbf{Answer:} B
\\ \textit{Options:} A) 6; B) 32; C) 30; D) 12; E) 36
\\ \textit{Note:} Correct option: B - 32
\item \textbf{Answer:} E
\\ \textit{Options:} A) conducting the survey at one shoe store; B) conducting the survey online targeting only the mall's registered customers; C) conducting the survey at all clothing stores; D) conducting the survey only on weekends; E) conducting the survey at the entrance to the mall
\\ \textit{Note:} Correct option: E - conducting the survey at the entrance to the mall
\item \textbf{Answer:} B
\\ \textit{Options:} A) 85.0; B) 60.0; C) 65.0; D) 45.0; E) 55.0
\\ \textit{Note:} Correct option: B - 60.0
\item \textbf{Answer:} A
\\ \textit{Options:} A) 13 in.; B) 10 in.; C) 11 in.; D) 45.0 in.; E) 15 in.
\\ \textit{Note:} Correct option: A - 13 in.
\item \textbf{Answer:} A
\\ \textit{Options:} A) 3.5; B) 1; C) 0.25; D) 0.5; E) 2.5
\\ \textit{Note:} Correct option: A - 3.5
\end{enumerate}

\section*{True / False}
\begin{enumerate}[label=\arabic*.]
\setcounter{enumi}{5}
\item \textbf{Answer:} True
\\ \textit{Options:} A) True; B) False
\\ \textit{Note:} LLM-generated true statement based on MCQ.
\item \textbf{Answer:} True
\\ \textit{Options:} A) True; B) False
\\ \textit{Note:} LLM-generated true statement based on MCQ.
\item \textbf{Answer:} True
\\ \textit{Options:} A) True; B) False
\\ \textit{Note:} LLM-generated true statement based on MCQ.
\item \textbf{Answer:} False
\\ \textit{Options:} A) True; B) False
\\ \textit{Note:} LLM-generated false statement. Correct answer: '6'.
\item \textbf{Answer:} False
\\ \textit{Options:} A) True; B) False
\\ \textit{Note:} LLM-generated false statement. Correct answer: '240'.
\end{enumerate}

\section*{Long Form Response}
\begin{enumerate}[label=\arabic*.]
\setcounter{enumi}{10}
\item \textbf{Answer:} We begin by factoring $336$. $336 = 2^4 \cdot 3 \cdot 7.$ Because we are looking for phone numbers, we want four single digits that will multiply to equal $336.$ Notice that $7$ cannot be multiplied by anything, because $7 \cdot 2$ is $14,$ which is already two digits. So, one of our digits is necessarily $7.$ The factor $3$ can be multiplied by at most $2,$ and the highest power of $2$ that we can have is $2^3 = 8.$ Using these observations, it is fairly simple to come up with the following list of groups of digits whose product is $336:$ \begin{align*} \&1, 6, 7, 8\\ \&2, 4, 6, 7\\ \&2, 3, 7, 8 \\ \&3, 4, 4, 7 \end{align*}For the first three groups, there are $4! = 24$ possible rearrangements of the digits. For the last group, $4$ is repeated twice, so we must divide by $2$ to avoid overcounting, so there are $\frac{4!}{2} = 12$ possible rearrangements of the digits. Thus, there are $3 \cdot 24 + 12 = \boxed{84}$ possible phone numbers that can be constructed to have this property.
\\ \textit{Note:} Students should show all steps in their mathematical solution.
\item \textbf{Answer:} Since $\angle ABP=90^{\circ}$, $\triangle ABP$ is a right-angled triangle. By the Pythagorean Theorem, $BP^2=AP^2-AB^2=20^2-16^2=144$$ and so $BP=12$, since $BP\textgreater{}0$. Since $\ensuremath{\angle} QTP=90^\{\ensuremath{\circ}\}$, $\ensuremath{\triangle} QTP$ is a right-angled triangle with $PT=12$. Since $PT=BP=12$, then by the Pythagorean Theorem, $$QT^2=QP^2-PT^2=15^2-12^2 = 81$$ and so $QT=9$, since $QT\textgreater{}0$. Our final answer is then $\ensuremath{\boxed{12,9}}\$.
\\ \textit{Note:} Students should show all steps in their mathematical solution.
\end{enumerate}

\vfill
\noindent\textit{End of Answer Key}
\end{document}